\originalchapter{Bose--Einstein凝聚}\label{ch:bose}
\originalsection{激发态容量与基态积聚}
自由Boson设基态能量为0. 巨正则$\mu<0$, 随降温在固定总数下$\mu$趋$0^-$. 将基态单独取出, 用$D(\epsilon)\propto\sqrt\epsilon$积分激发态, 记$z=\e^{\beta\mu}\le1$, 定义$g_s(z)=\sum_{r\ge1}z^r/r^s$. 展开几何级数并积分$\int\epsilon^{1/2}\e^{-r\beta\epsilon}\dd\epsilon$, 得
\[
N_{\rm ex}=\frac{gV}{\lambda_T^3}g_{3/2}(z),\qquad
g_{3/2}(1)=\zeta(3/2)\simeq2.612375.
\]
激发态在给定温度最多容纳$gV\zeta(3/2)/\lambda_T^3$个粒子. 自旋0且固定密度$n$的临界温度与凝聚比例为
\[
T_c=\frac{2\pi\hbar^2}{m\kb}\left(\frac n{\zeta(3/2)}\right)^{2/3},\qquad
\frac{N_0}N=1-(T/T_c)^{3/2}\quad(T\le T_c).
\]
超过激发态容量的粒子进入单一基态, 这就是宏观占据. 不可将基态贡献藏进连续态密度积分的单点; 连续积分的单点测度为零, 而$N_0$是量级$N$的独立项.

固定有限$N,V$时$\mu$仍稍低于0, $N_0=1/(\e^{-\beta\mu}-1)$; $\mu=0$是热力学极限的简写. 理想巨正则基态有量级$N_0^2$的数涨落, 若实际系统固定总数, 凝聚态涨落需用相应系综, 不能把这个“巨正则灾难”当成实验必然结论.

\originalsection{能量、热容与空间维数}
激发态能量积分给$U=(3/2)gV\kb T\lambda_T^{-3}g_{5/2}(z)$和$P=2U/(3V)$. 低于$T_c$取$z=1$, 基态能量0, 故$U\propto T^{5/2}$,
\[
\frac{C_V}{N\kb}=\frac{15}{4}\frac{\zeta(5/2)}{\zeta(3/2)}(T/T_c)^{3/2}.
\]
高于$T_c$, $N$守恒要求$\dd\ln z/\dd\ln T=-(3/2)g_{3/2}/g_{1/2}$, 因而
\[
\frac{C_V}{N\kb}=\frac{15}{4}\frac{g_{5/2}(z)}{g_{3/2}(z)}
-\frac94\frac{g_{3/2}(z)}{g_{1/2}(z)}.
\]
在$z\to1^-$时$g_{1/2}$发散, 热容两侧连续但导数出现尖点; 高温$z\ll1$恢复$3/2$. 在均匀理想$d$维气体中, 低能态密度$D(\epsilon)\propto\epsilon^{d/2-1}$, Bose占据约$\kb T/\epsilon$, 积分近零收敛需$d>2$. 二维和一维均匀理想气体在无限体积没有有限温度的这种凝聚, 但陷阱或相互作用改变低能结构, 必须另作讨论.

\originalsection{陷阱与观察}
三维谐振阱$V=\frac12m(\omega_x^2x^2+\omega_y^2y^2+\omega_z^2z^2)$, 在$\kb T$大于级间距的连续近似中, 激发态累计数是能量单纯形的体积, 给$D(\epsilon)=\epsilon^2/(2\hbar^3\omega_x\omega_y\omega_z)$. 因而$N_{ex}=\zeta(3)(\kb T)^3/(\hbar^3\omega_x\omega_y\omega_z)$, $N_0/N=1-(T/T_c)^3$. 它与均匀盒的$3/2$幂不同, 不能把一种实验的凝聚比例用于另一种容器.

激光冷却可以降低特定原子速度, 磁阱避免原子接触壁面; 蒸发冷却优先放走高能原子, 剩余粒子经碰撞再热化, 使温度下降. 降温同时损失粒子, 是否接近凝聚要看相空间密度$n\lambda_T^3$, 不能只看温度计读数. 自旋极化氢避免快速复合, 但壁面自旋翻转使早期壁囚禁方案困难; 磁囚禁和蒸发改变了这一限制.

释放陷阱后, 热云以热速度扩展, 基态凝聚的动量宽度由量子尺寸决定, 形成窄峰与宽背景的双组分分布. 在理想抛物阱中热云空间与动量尺度随$\sqrt T$变化, 基态尺度为$\sqrt{\hbar/(m\omega)}$和$\sqrt{m\hbar\omega}$. 相互作用会改变膨胀动力学, 本稿的理论示意不替代原课实验照片. 两凝聚云重叠可以观察干涉条纹, 展示宏观物质波的相干性.

Bose与Einstein于1924--1925年提出理想气体凝聚. 原课的历史线索还包括液氦低温技术、超导、${}^4$He及${}^3$He超流, 到1995年稀薄碱金属气体凝聚及其后氢凝聚. 超导的载流配对与超流的无耗散响应都涉及相互作用; BEC是单粒子态宏观占据, 三者有关却不等价. 本册从这些现象引向相互作用统计, 不把历史图片或奖项引文当作物理机制的证明.
